% Chương 3. Phụ trách: Minh (MINH-8)
% Nguồn: docs/model_spec.md phần A và B.

\chapter{Mô hình hóa bằng quy hoạch nguyên tuyến tính}
\label{ch:mo-hinh}

\section{Tổng quan ILP và lý do lựa chọn}
\label{sec:ilp}

\section{Bài toán A: lịch thi}
\label{sec:mo-hinh-a}

\subsection{Tập hợp, tham số và biến quyết định}
\label{subsec:a-bien}
% model_spec A.2-A.4

\subsection{Ràng buộc cứng}
\label{subsec:a-cung}
% model_spec A.5

\subsection{Ràng buộc mềm và hàm mục tiêu}
\label{subsec:a-mem}
% model_spec A.6

\subsection{Chiến lược giải một pha và hai pha}
\label{subsec:a-chien-luoc}
% model_spec A.7

\subsection{Kích thước mô hình}
\label{subsec:a-kich-thuoc}
% model_spec A.11

\section{Bài toán B: lịch phòng học và phòng máy}
\label{sec:mo-hinh-b}

\subsection{Tập hợp, tham số và biến quyết định}
\label{subsec:b-bien}
% model_spec B.2-B.4

\subsection{Ràng buộc cứng}
\label{subsec:b-cung}
% model_spec B.5

\subsection{Ràng buộc mềm và hàm mục tiêu}
\label{subsec:b-mem}
% model_spec B.6

\subsection{Chiến lược giải và hậu xử lý}
\label{subsec:b-chien-luoc}
% model_spec B.8

\section{Xếp lại khi có sự cố}
\label{sec:xep-lai}
% model_spec A.14 và B.9

\section{Kiểm tra dữ liệu và chẩn đoán khi vô nghiệm}
\label{sec:chan-doan}
% model_spec A.8 và B.7

\section{Đối chiếu với 10 quy tắc lập mô hình của Daskin}
\label{sec:daskin}
% model_spec A.9 và B.10

\section{Ví dụ toy tính tay}
\label{sec:vi-du-toy}
% model_spec A.13, B.12; data/toy, data/toy_b

\section{Những điều mô hình cố ý không xét}
\label{sec:khong-xet}
% model_spec A.12 và B.11
